% ============================================================
%  Offline Training with Pre-Gathered Human Feedback
%  Adaptation of main_oracle.py for post-study aggregation
% ============================================================
\documentclass[11pt, a4paper]{article}

\usepackage[margin=2.5cm]{geometry}
\usepackage{amsmath, amssymb}
\usepackage{algorithm}
\usepackage{algpseudocode}
\usepackage{booktabs}
\usepackage{hyperref}
\usepackage{xcolor}
\usepackage{listings}
\usepackage{parskip}
\usepackage{graphicx}
\usepackage{caption}
\usepackage{subcaption}

\hypersetup{colorlinks=true, linkcolor=blue, citecolor=blue, urlcolor=blue}

\lstset{
  basicstyle=\ttfamily\small,
  keywordstyle=\color{blue},
  commentstyle=\color{gray},
  backgroundcolor=\color{gray!8},
  frame=single,
  breaklines=true,
  language=Python,
}

\title{\textbf{Offline Training with Pre-Gathered Human Feedback}\\
  \large Adapting the Oracle Training Loop for Post-Study Aggregation}
\author{RL-Multiple-Human-Feedback Pipeline}
\date{\today}

% ============================================================
\begin{document}
\maketitle
\tableofcontents
\newpage

% ============================================================
\section{Background and Motivation}
\label{sec:background}

The system uses \texttt{tabQL\_Cest\_vi\_t2}: tabular Q-learning augmented with
policy shaping from human feedback, where each trainer's consistency is
estimated via Variational Inference (VI).  During a live study, human
participants label Pac-Man game states through a web interface; each label
immediately updates the agent's brain stored in Google Cloud Storage (GCS).
Because each participant runs an independent session (\texttt{nTrainer}$=1$),
their feedback is never aggregated until after the study ends.

The \emph{post-study aggregation script} (\texttt{offline\_train.py}) was
therefore needed to:
\begin{enumerate}
  \item Download all participant brains from GCS.
  \item Stack their feedback arrays into a single multi-trainer agent
        (\texttt{nTrainer}$=N$).
  \item Estimate per-participant consistency weights $C_e[m]$.
  \item Train and compare: \emph{combined feedback}, \emph{individual
        participant feedback}, and \emph{pure RL baseline}.
\end{enumerate}

Performance on the offline script was found to be substantially lower than
on the oracle benchmark (\texttt{main\_oracle.py}), even with equivalent
episode counts.  This document describes the root causes and the changes made
to close that gap.

% ============================================================
\section{The Oracle Training Procedure}
\label{sec:oracle}

\texttt{main\_oracle.py} provides the performance reference.  It trains an
agent online, generating synthetic human feedback at every environment step
from a pre-trained greedy oracle policy.  The oracle's Q-table encodes a
near-optimal policy; at each step it compares the oracle's preferred action
with the agent's current state and produces a \texttt{Feedback} object
accordingly.

\subsection{Key parameters}
\begin{table}[h!]
\centering
\begin{tabular}{lll}
\toprule
Parameter & Value & Description \\
\midrule
\texttt{episode\_count}       & 1000        & Training episodes per trial \\
\texttt{max\_steps}           & 500         & Maximum steps per episode \\
\texttt{reset\_env\_random}   & \texttt{False} & Fixed start state every episode \\
\texttt{update\_Cest\_interval} & 5         & Re-estimate $C_e$ every $N$ episodes \\
\texttt{active\_feedback\_type} & \texttt{count} & Active learning: label if $N_{sa} < \theta$ \\
\texttt{feedback\_type}       & \texttt{ranked} & Feedback format \\
\bottomrule
\end{tabular}
\caption{Oracle training parameters.}
\end{table}

\subsection{Per-step feedback injection}

At each environment step $j$ within episode $i$, the oracle generates a
\texttt{Feedback} object for the current trajectory:

\begin{algorithm}[H]
\caption{Oracle training loop (simplified from \texttt{main\_oracle.py})}
\begin{algorithmic}[1]
\For{$i = 1, \ldots, \texttt{episode\_count}$}
  \State \Call{reset\_episode}{random=False}
  \For{$j = 1, \ldots, \texttt{max\_steps}$}
    \State $fb \gets$ \Call{generate\_feedback}{trajectory, oracle, active\_learning}
    \State $a, s', r, done \gets$ \Call{trainer.step}{feedback=$fb$, update\_Cest=False}
    \If{$done$ \textbf{or} $j = \texttt{max\_steps}$}
      \State $fb_{\text{end}} \gets$ \Call{generate\_feedback}{\ldots, end\_of\_episode=True}
      \State update\_Cest $\gets (i \bmod \texttt{interval} = 0)$
      \State \Call{agent.act}{$a,\ s',\ r,\ done,\ fb_{\text{end}},\ C,$\ update\_Cest}
      \State $\texttt{prev\_obs} \gets \texttt{None}$\quad\textit{// reset for next episode}
      \State \textbf{break}
    \EndIf
  \EndFor
\EndFor
\end{algorithmic}
\end{algorithm}

Crucially, injecting $fb$ on every step means that whenever the agent visits a
state that the oracle would label, the feedback arrays $h^+$ and $h^-$
\emph{accumulate} throughout training.  This has two downstream effects:
\begin{itemize}
  \item The digamma shaping values $\psi^r[m]$ and $\psi^w[m]$ strengthen
        over time as more feedback accumulates.
  \item $C_e[m]$ rises from the initial $0.5$ towards its asymptote as the EM
        loop has more evidence to work with.
\end{itemize}

% ============================================================
\section{Deficiencies in the Original Offline Script}
\label{sec:deficiencies}

The original \texttt{offline\_train.py} differed from the oracle procedure in
three important ways.

\subsection{Frozen feedback arrays (most significant)}

The original offline script loaded $h^+$ and $h^-$ from GCS, ran one
C-estimation pass to initialise $\psi^r$ and $\psi^w$, and then trained with
\emph{empty} feedback at every step:

\begin{lstlisting}
empty_fb = [[] for _ in range(N)]
while not trainer.done:
    trainer.step(feedback=empty_fb, ...)   # hp/hm never grow
\end{lstlisting}

Because \texttt{\_collect\_feedback} is only called when
\texttt{prev\_obs is not None} \emph{and} the feedback list is non-empty, the
human feedback arrays remained frozen at their study-time values.  The shaping
signal was therefore at its weakest at the start of training and never
strengthened, unlike the oracle where feedback accumulates continuously.

\subsection{Missing terminal Q-update}

The original loop exited when \texttt{trainer.done} became \texttt{True}, then
reset \texttt{prev\_obs = None} \emph{before} any explicit terminal update:

\begin{lstlisting}
while not trainer.done:
    trainer.step(...)                   # last step sets done=True, exits loop
trainer.agent.prev_obs = None           # blocks Q-update
trainer.agent.act(..., update_Cest=True)  # Ce only, no Q update
\end{lstlisting}

Because the Q-update inside \texttt{tabQL\_Cest\_vi} is guarded by
\texttt{if self.prev\_obs is not None}, the terminal transition reward was
never propagated into $Q$.  In contrast, the oracle explicitly calls
\texttt{agent.act()} once more after the terminal step (line 112 of
\texttt{main\_oracle.py}), before resetting \texttt{prev\_obs}.

\subsection{Random episode start states}

The oracle uses \texttt{reset\_env\_random=False}, placing the agent at the
same start position every episode.  This ensures consistent trajectories
through the labelled region of the state space, which is critical when
feedback coverage is sparse.  The original offline script called
\texttt{reset\_episode()} with the default \texttt{random=True}, causing
the agent to start from arbitrary positions and frequently miss states that
were labelled during the study.

% ============================================================
\section{Changes Made to \texttt{offline\_train.py}}
\label{sec:changes}

Two training modes were introduced, selectable via \texttt{--training-mode}.

\subsection{Replay mode (default) — mirrors \texttt{main\_oracle.py}}

\paragraph{Core idea.}
Rather than passing empty feedback every step, the frozen study arrays
$h^+_0, h^-_0$ (a copy taken before training begins) are used as a
\emph{lookup table}.  Whenever the agent visits a state that was labelled
during the study, and active learning indicates that state still needs
coverage, a \texttt{Feedback} object is synthesised from $h^+_0 - h^-_0$
and injected — identically to how the oracle would inject feedback.

\paragraph{Active-learning gate.}
The count-based active-learning criterion from \texttt{main\_oracle.py} is
preserved: feedback is only injected for state $s$ if
$N_{sa}[s, a] < \theta_{\text{AL}}$ for any action $a$.  This limits how many
times any state receives additional label mass, preventing the human signal
from being inflated relative to the oracle's single-label-per-visit
behaviour.  The default threshold $\theta_{\text{AL}} = 3$ matches the
oracle's active-learning configuration.

\paragraph{Feedback construction.}
For each trainer $m$ and current state $s$, the net signal is:
\begin{equation}
  \Delta^m_a = h^{+,m}_{0,s,a} - h^{-,m}_{0,s,a}, \quad a \in \mathcal{A}
\end{equation}
Actions with $\Delta^m_a > 0$ are placed in \texttt{good\_actions}; those
with $\Delta^m_a < 0$ in \texttt{bad\_actions}.  This is analogous to the
oracle comparing its Q-table against the agent's and labelling accordingly.
The \texttt{Feedback} object is passed to \texttt{trainer.step()}, which
calls \texttt{\_collect\_feedback(fb)}, incrementing the live $h^+$ and
$h^-$ arrays.

\paragraph{Terminal Q-update.}
After the episode-ending step the script now calls \texttt{agent.act()}
explicitly — before resetting \texttt{prev\_obs} — so that the terminal
reward is propagated into $Q$:

\begin{lstlisting}
if done or steps >= max_steps:
    update_Cest = ((ep + 1) % update_Cest_interval == 0)
    trainer.agent.act(action_idx, ob, rw, done,
                      fb, 0.5, update_Cest=update_Cest)
    trainer.agent.prev_obs = None   # reset AFTER Q-update
    break
\end{lstlisting}

\paragraph{Fixed start state.}
\texttt{trainer.reset\_episode(random=False)} is used throughout the replay
training loop, matching the oracle's \texttt{reset\_env\_random=False}.

\paragraph{C-estimation frequency.}
$C_e$ is re-estimated every \texttt{--cest-interval} episodes (default 5),
matching the oracle's \texttt{update\_Cest\_interval=5}.

The full replay loop is given in Algorithm~\ref{alg:replay}.

\begin{algorithm}[H]
\caption{Replay training loop (\texttt{train\_with\_human\_replay})}
\label{alg:replay}
\begin{algorithmic}[1]
\Require Frozen study arrays $h^+_0, h^-_0$; trainer with loaded $h^+, h^-$
\For{$ep = 1, \ldots, \texttt{n\_episodes}$}
  \State \Call{reset\_episode}{random=\texttt{False}}
  \For{$j = 1, \ldots, \texttt{max\_steps}$}
    \State $s \gets \texttt{trainer.ob}$
    \If{$\min_a N_{sa}[s,a] < \theta_{\text{AL}}$
        \textbf{and} $\sum_{a,m} (h^{+,m}_{0,s,a} + h^{-,m}_{0,s,a}) > 0$}
      \State Build $fb$ from $\{h^+_{0,s,:} - h^-_{0,s,:}\}$ per trainer
    \Else
      \State $fb \gets [[\,]\,]_{m=1}^{N}$
    \EndIf
    \State $a, s', r, done \gets$ \Call{trainer.step}{feedback=$fb$, update\_Cest=False}
    \If{$done$ \textbf{or} $j = \texttt{max\_steps}$}
      \State $\text{updateCe} \gets (ep \bmod \texttt{interval} = 0)$
      \State \Call{agent.act}{$a,\ s',\ r,\ done,\ fb,\ 0.5,$\ updateCe}
      \State $\texttt{prev\_obs} \gets \texttt{None}$
      \State \textbf{break}
    \EndIf
  \EndFor
\EndFor
\end{algorithmic}
\end{algorithm}

\subsection{Static mode — original behaviour with bug fix}

The static mode retains the original approach (frozen $h^+/h^-$, fixed psi
values) but incorporates the terminal Q-update fix and the
\texttt{update\_Cest\_interval} parameter.  It is available via
\texttt{--training-mode static} and is useful as a controlled comparison to
quantify the benefit of the replay mechanism independently of other changes.

% ============================================================
\section{Algorithmic Detail: Policy Shaping}
\label{sec:policy_shaping}

For completeness, the policy shaping mechanism in \texttt{tabQL\_Cest\_vi\_t2}
(type~2, ``only one optimal action'') is reproduced here.  At each step the
Boltzmann log-probability over actions is modified by the accumulated human
feedback:

\begin{equation}
  \ell_a \;=\; \frac{Q(s,a)}{\tau}
  \;+\; \sum_{m=1}^{N}
    \Bigl[
      h^{+,m}_{s,a}\,\psi^r_m
      + h^{-,m}_{s,a}\,\psi^w_m
      - h^{-,m}_{s,a}\,\psi^r_m
      - h^{+,m}_{s,a}\,\psi^w_m
    \Bigr]
\end{equation}

where $\tau = 1.5$ is the temperature constant, and the digamma shaping
values for trainer $m$ are:

\begin{align}
  \psi^r_m &= \psi\!\left(\textstyle\sum h^r_m + \alpha_m\right)
             - \psi\!\left(\textstyle\sum h^r_m + \sum h^w_m + \alpha_m + \beta_m\right)
  \\[4pt]
  \psi^w_m &= \psi\!\left(\textstyle\sum h^w_m + \beta_m\right)
             - \psi\!\left(\textstyle\sum h^r_m + \sum h^w_m + \alpha_m + \beta_m\right)
\end{align}

Here $\psi(\cdot)$ is the digamma function, $\alpha_m, \beta_m$ are the Beta
prior parameters (\texttt{prior\_alpha}, \texttt{prior\_beta}; default $1.0$
for a flat prior), and $\sum h^r_m,\,\sum h^w_m$ are the running totals
updated by the EM loop during C-estimation.

The key consequence for the two training modes is:
\begin{itemize}
  \item \textbf{Static mode}: $h^+, h^-$ are frozen so $\psi^r_m, \psi^w_m$
        change only when C-estimation runs, driven solely by the Q-table
        evolving relative to fixed feedback counts.
  \item \textbf{Replay mode}: $h^+, h^-$ grow each time the agent revisits a
        labelled state and the active-learning gate passes.  Both $\psi$ values
        therefore strengthen monotonically, producing an increasingly confident
        shaping signal as training progresses — precisely mirroring the
        oracle's behaviour.
\end{itemize}

% ============================================================
\section{Differences Remaining vs.\ the Oracle}
\label{sec:remaining_differences}

Even with replay mode, one structural difference from the oracle remains:

\paragraph{Fixed vs.\ continuously generated labels.}
The oracle generates a fresh label for \emph{every} state it visits, based on
its current Q-table vs.\ the agent's Q-table.  The replay script can only
label states that were seen during the study.  States never labelled by any
participant contribute no shaping signal regardless of the active-learning
threshold.  In practice this means:
\begin{itemize}
  \item Coverage is bounded by what participants happened to see.
  \item States visited during training but absent from $h^+_0/h^-_0$ receive
        no reinforcement from human feedback, so Q-learning alone drives
        behaviour there.
\end{itemize}
This is an inherent limitation of offline aggregation and motivates ensuring
broad state coverage during the study (e.g.\ larger bundle sizes, more
participants, or entropy-based active learning during data collection).

% ============================================================
\section{New Command-Line Interface}
\label{sec:cli}

\begin{table}[h!]
\centering
\begin{tabular}{llp{7.5cm}}
\toprule
Flag & Default & Description \\
\midrule
\texttt{--training-mode} & \texttt{replay}
  & \texttt{replay}: inject labels at each step (oracle-equivalent).
    \texttt{static}: frozen $h^+/h^-$ shaping prior. \\
\texttt{--al-threshold}  & \texttt{3.0}
  & Active-learning visit count threshold; feedback injected while
    $N_{sa} < \theta$. \\
\texttt{--cest-interval} & \texttt{5}
  & Re-estimate $C_e$ every $N$ episodes (matches oracle default). \\
\texttt{--n-episodes}    & \texttt{100}
  & Training episodes per condition.  Recommend $\geq 1000$ for
    convergence at oracle scale. \\
\texttt{--q-init}        & \texttt{zeros}
  & \texttt{zeros}: fair comparison baseline.
    \texttt{average}: warm-start from participants' Q-tables. \\
\texttt{--sessions}      & all
  & Participant codes to include; omit to use all brains in bucket. \\
\bottomrule
\end{tabular}
\caption{Updated \texttt{offline\_train.py} CLI flags.}
\end{table}

\paragraph{Recommended invocation (oracle-equivalent).}
\begin{lstlisting}[language=bash]
python offline_train.py \
    --bucket jerskine_human_feedback \
    --sessions P001 P002 P003 \
    --n-episodes 1000 \
    --training-mode replay          # default, shown for clarity
\end{lstlisting}

\paragraph{Comparison run.}
\begin{lstlisting}[language=bash]
# Replay mode
python offline_train.py ... --training-mode replay --save results_replay.pkl
# Static mode (original behaviour, with terminal bug fix)
python offline_train.py ... --training-mode static --save results_static.pkl
# Plot both
python analyse_results.py results_replay.pkl --save plot_replay.png
python analyse_results.py results_static.pkl --save plot_static.png
\end{lstlisting}

% ============================================================
\section{Summary of Changes}
\label{sec:summary}

\begin{table}[h!]
\centering
\begin{tabular}{p{4cm} p{4.5cm} p{4.5cm}}
\toprule
Aspect & Before & After (replay mode) \\
\midrule
Feedback during training
  & Empty \texttt{[[]}]} at every step; $h^+/h^-$ frozen
  & \texttt{Feedback} object injected for labelled states passing
    active-learning gate; $h^+/h^-$ accumulate \\[4pt]
Terminal Q-update
  & \texttt{prev\_obs=None} before terminal \texttt{act()};
    terminal reward never hits $Q$
  & Explicit \texttt{act(action, s', r, done=True, ...)} before
    \texttt{prev\_obs=None} \\[4pt]
Episode start state
  & \texttt{random=True} (arbitrary)
  & \texttt{random=False} (fixed, matches oracle) \\[4pt]
$C_e$ update frequency
  & Every episode
  & Every \texttt{--cest-interval} episodes (default 5) \\[4pt]
$\psi^r / \psi^w$ over training
  & Constant after initial C-estimation
  & Grow as $h^+/h^-$ accumulate \\[4pt]
Episode count
  & Default 100
  & Recommend $\geq 1000$ \\
\bottomrule
\end{tabular}
\caption{Comparison of original offline training vs.\ updated replay mode.}
\end{table}

\end{document}
